\section{Sprint 2 -- Preliminary Design Review (PDR)}
%=============================================================================

\begin{priorities}
\begin{enumerate}[leftmargin=*, itemsep=1pt, topsep=0pt]
    \item Thoroughly reverse engineer the COTS product, especially its electrical interfaces.
    \item Complete an honest FMEA; do not downplay severity to look ``clean.''
    \item Weight your decision matrix criteria \textit{before} scoring alternatives.
\end{enumerate}
\end{priorities}

\begin{checklistbox}[Before You Start: Sprint 2 Preparation]
\begin{enumerate}[leftmargin=*, itemsep=1pt, topsep=0pt]
    \item Read MRD Chapters~4 (Concept Generation) and 21 (PDR deliverables).
    \item Review your Sprint~1 requirements. Are they all testable and quantified?
    \item Obtain the COTS product and identify all interfaces (mechanical, electrical, data).
    \item Prepare functional decomposition before selecting any components.
    \item Gather materials for the morphological chart (at least 3 options per function).
\end{enumerate}
\end{checklistbox}

\begin{protip}
\textbf{Start with functional decomposition before jumping to components.} Before picking sensors, microcontrollers, or actuators, list the \textit{functions} your system must perform using verb-noun pairs: Sense Temperature, Dispense Food, Filter Air, etc. Then break each function into sub-functions. This separates \textit{what} the system must do from \textit{how} it does it, and it directly feeds your trade study. A morphological chart (listing functions as rows and 3--5 implementation options as columns) is the most powerful way to systematically generate meaningfully different design concepts. See MRD Chapters~3--4 (Requirements \& Concept Generation); follow the teal GreenShield boxes for worked examples.
\end{protip}

\subsection{Reverse Engineering}

\begin{faqbox}[Q: What am I actually looking for when I reverse engineer the COTS product?]
You are trying to answer: ``What constraints does this existing product impose on my add-on design?'' Look for mounting points, electrical interfaces (voltage, current, connector types), physical dimensions and tolerances, material properties, and any features that could either help or hinder your integration. Document everything with photos and measurements.
\end{faqbox}

\begin{protip}
Measure twice (or three times). Tolerances on consumer products can be surprisingly loose. If you are designing a mechanical interface, build in adjustability or compliance rather than relying on a single measurement from one unit.
\end{protip}

\begin{warnbox}
Explicitly document the \textbf{connectors and power draw} of the COTS product during reverse engineering. One of the most common integration failures occurs when an add-on system pulls more current than the parent product's power supply can handle, or when the team orders a mating connector with the wrong pitch (e.g., 2.0\,mm vs.\ 2.54\,mm JST headers). Record connector type, pin count, pitch, voltage, and rated/measured current draw. For any stainless steel fluid fittings, apply nickel anti-seize compound before assembly to prevent thread galling (a cold-weld that permanently destroys both parts). This information is critical for your CDR and will prevent expensive, time-consuming mistakes in Sprints~3--4.
\end{warnbox}

\subsection{Failure Modes and Effects Analysis (FMEA)}

\begin{tipbox}
\textbf{Start your FMEA now, not at FDR.} The Master Reference Document positions FMEA as a design-phase tool (Chapters~2--3), not a post-mortem. Performing FMEA during architecture enables you to design out critical failure modes before ordering parts. Waiting until the final report turns FMEA into an autopsy.
\end{tipbox}

\begin{faqbox}[Q: How do I identify failure modes I have not thought of yet?]
Systematically walk through every component and every interface in your system and ask: ``What happens if this breaks, jams, shorts, overheats, receives the wrong input, or loses power?'' Also consider misuse scenarios: what if the user does something unexpected? The USCSB ExxonMobil video assigned in this sprint is an excellent example of how overlooked failure modes cascade into catastrophic outcomes.
\end{faqbox}

\begin{warnbox}
Do not assign uniformly low severity or occurrence ratings to make your FMEA look ``clean.'' An FMEA with no high-risk items is almost certainly an FMEA that was not done thoroughly. Instructors will see through this immediately. The value of the FMEA is in \textit{finding} the scary stuff so you can address it in your design.
\end{warnbox}

\begin{tipbox}
\textbf{Score the FMEA using three independent scales.} Each failure mode gets rated on Severity (how bad if it happens, 1--10), Occurrence (how likely, 1--10), and Detection (how likely the failure goes unnoticed, 1--10). Multiply them to get the Risk Priority Number: $\text{RPN} = S \times O \times D$. Failure modes with RPN $\geq 100$ generally need corrective action in your design. A high-severity, low-occurrence, hard-to-detect failure (e.g., temperature sensor reading high so frost is not detected ($S=9, O=3, D=5, \text{RPN}=135$), demands more attention than an obvious, low-impact failure. Prioritize your design resources based on RPN, not gut feeling.
\end{tipbox}

\begin{protip}
\textbf{Think about all six types of interfaces during your FMEA.} Failures cluster at interfaces, not inside components. The six interface categories are: electrical (voltage, protocol, grounding), mechanical (mounting, tolerances, alignment), thermal (heat paths, contact resistance, component overheating), software/data (APIs, data format, timing), fluid (fittings, pressure, seal compatibility), and human (display visibility, control layout, error prevention). Walk through each interface in your system and ask what happens when each one fails.
\end{protip}

\subsection{PDR Report and Decision Matrix}

\begin{faqbox}[Q: How many design alternatives should we consider in the trade study?]
At minimum, you should have three meaningfully different concepts. ``Meaningfully different'' means different architectures or approaches, not minor variations of the same idea (e.g., ``same design but in blue'' does not count). The decision matrix should use weighted criteria that trace back to your customer needs and SDRD objectives.
\end{faqbox}

\begin{tipbox}
When building your decision matrix, weight the criteria \textit{before} scoring the alternatives. This prevents the common bias of adjusting weights after the fact to justify the concept you already wanted to build.
\end{tipbox}

\begin{protip}
\textbf{Run a sensitivity analysis on your decision matrix.} After scoring, vary each weight by $\pm$0.05 and check whether the winning concept changes. For instance, if Concept~A beats Concept~B by only 0.05 points, and shifting the ``cost'' weight by 10\% makes Concept~B the winner, your decision is not robust; you need better data on that criterion or should prototype both options. A margin of less than 5\% between the top two concepts demands additional investigation, not a hasty commitment. Also ensure every candidate is a legitimate contender: including a deliberately weak ``straw man'' option to make your preferred concept win is a common pitfall that instructors spot immediately.
\end{protip}

\begin{faqbox}[Q: How do I know if my PDR package is ready?]
Apply the \textbf{PDR Litmus Test}: ``If a different engineering team read only our PDR package, could they independently arrive at the same concept selection and understand why it is the best choice?'' If the answer is yes, your PDR is defensible. If the answer requires ``well, they would also need to know that\ldots,'' that missing information must be added to the package. Remember, PDR answers the question ``\textit{Should} this design work?'' Your deliverables must demonstrate that the problem is well-defined, the requirements are sound, and the selected concept is feasible.
\end{faqbox}

%=============================================================================
%=============================================================================
\vspace{0.5em}
\begin{center}
\textcolor{minesblue}{\rule{0.6\textwidth}{1pt}}\\[4pt]
{\Large\bfseries\sffamily\color{minesblue} Phase 2: ``Will This Design Work?''}\\[2pt]
{\small\itshape Sprints 3--4 freeze the design, procure parts, and demonstrate critical subsystems.}\\[4pt]
\textcolor{minesblue}{\rule{0.6\textwidth}{1pt}}
\end{center}
\vspace{0.5em}
%=============================================================================
